\documentclass[twocolumn]{preport}

\usepackage[dvipdfmx]{graphicx}

\graphicspath{{figs/}}

\title{1192年度 研究要旨:\\
レポートを執筆する環境に関する研究}
\author{1192 源頼朝}

\begin{document}
\pagestyle{empty}
\maketitle
\thispagestyle{empty}
\sloppy

\section{はじめに}

本稿はレポートのテンプレートである．

まずはじめに参考文献を引用する\cite{Knuth86}．

さらにもうひとつ引用する\cite{BollenHHJ07}．

日本語の文献も引用する\cite{Okumura07}．

最後にウェブページを引用する\cite{Github}．

\section{表や図の描画方法}

図は\figref{sample}，表は\tabref{sample}として参照する．

\begin{figure}[tbh]
 \centering
 \begin{minipage}{0.3\columnwidth}
  \includegraphics[width=\columnwidth]{nowprinting.eps}
  \caption{eps図の参考例}
 \end{minipage}
 \hspace{0.15\columnwidth}
 \begin{minipage}{0.3\columnwidth}
  \includegraphics[width=\columnwidth]{dj.jpg}
  \caption{jpg図の参考例}
 \end{minipage}
 \figlab{sample}
\end{figure}

\begin{table}[tbh]
 \centering
 \begin{tabular}{|l|r|} \hline
 A1 & B1 \\
 A2 & B2 \\ \hline
 \end{tabular}
 \caption{図の参考例}
 \tablab{sample}
\end{table}

\section{おわりに}

本稿において本文を通してレポートの執筆が可能な環境であることが示された．

\bibliographystyle{junsrt}
\bibliography{main}

\end{document}